\scanpage{127}
\mspage{1}
% archive sheet number: 126
\uline{Théorèmes} de \uline{factorisation dans les catégories accessibles}

% slanted margin heading, read after rotation
\margin{1. \uline{Construction} \uline{fondamentale} : la suite transfinie de factorisations d'une flèche $f$}

Soit $M$ une \inserted{$\unsure{\mathcal U|U}$-}catégorie, stable par ``petites'' (i.e. $\mathcal U$-petites) $\underset{\longrightarrow}{\lim}$ filtrantes. Soit
\[C_0\subset\mathrm{Fl}\,M\tag*{1.1.}\]
un \inserted{\uline{petit}} ensemble de flèches, on se propose de construire pour toute flèche $f$ de $M$, une factorisation
\[f=pi\qquad i\in\widetilde C_0=((C_0)_*)^{\unsure{\kappa|*}},\ p\in(C_0)_*.\tag*{1.2.}\]
Pour ceci, on va construire une suite ordinale de foncteurs
\[
\left\{\begin{aligned}
f&\longmapsto\deleted{\unsure{\Sigma|Z}_\alpha}\unsure{X_\alpha|\Sigma_\alpha|Z_\alpha}(f)=\deleted{\Sigma_\alpha}X_\alpha(C_0,f)\\
\deleted{\Sigma_\alpha}X_\alpha&:\underline{\mathrm{Fl}}(M)\longrightarrow\deleted{\underline{\mathrm{Fl}}(}M\deleted{)}
\end{aligned}\right.\tag*{1.3.}
\]
\deleted{\unsure{tels que|tel que}} et des hom. fonctoriels \deleted{\unsure{avec} $f:X\longrightarrow Y$}
\deleted{$\unsure{\operatorname{source}}\,\Sigma_\alpha(f)=$}
\diagram{\begin{tikzcd}[ampersand replacement=\&]
\deleted{X} \& \deleted{\unsure{\Sigma|X|Z}_\alpha(f)} \arrow[l,"\deleted{i_\alpha(f)}"']
\end{tikzcd}}
\deleted{$\operatorname{source}f\longrightarrow X_\alpha(f)$ \unsure{l'identité} \unsure{dont}}
\[i_\alpha:\operatorname{source}f\longrightarrow X_\alpha(f)\qquad\text{(fonctoriel en $f$)}\tag*{1.4.}\]
\deleted{et un hom. fonctoriel en $f$}
\diagram{\begin{tikzcd}[ampersand replacement=\&]
\deleted{\unsure{\operatorname{but}|\operatorname{source}}}\deleted{\Sigma_\alpha}X_\alpha(f) \arrow[r,"p_\alpha(f)"] \& \operatorname{but}(f)
\end{tikzcd}}
\[\tag*{1.5.}\]
de sorte qu'on a le diagramme (où $X=\operatorname{source}f$, $Y=\operatorname{but}f$)
\diagram{\begin{tikzcd}[ampersand replacement=\&]
X \arrow[rr,dash,"f"] \arrow[dr,"\deleted{\unsure{\Sigma|I}_\alpha}i_\alpha(f)"'] \& \& Y \\
\& X_\alpha(f) \arrow[ur,"p_\alpha(f)"'] \&
\end{tikzcd}}
\[\tag*{1.6.}\]
, ce diagramme étant commutatif
\[f=p_\alpha(f)\deleted{\Sigma_\alpha}i_\alpha(f).\tag*{1.7.}\]
On peut dire aussi qu'on a une section du foncteur canonique \inserted{\unsure{dit} ``\unsure{composition}''} (transposé de $\unsure{\mathbb{\Delta}|\Delta}^{1}\longrightarrow\mathbb{\Delta}^{2}$, donné par $0\mapsto0$, $1\mapsto2$ :
\[\tag*{1.8.}\]
\diagram{\begin{tikzcd}[ampersand replacement=\&,execute at end picture={\node[anchor=east] at ([xshift=-5pt]current bounding box.west) {$\Biggl\{$};}]
\underline{\operatorname{Nerf}}_2(M) \arrow[r,"\unsure{\sigma|\rho}"] \& \underline{\operatorname{Nerf}}_1(M)=\underline{\mathrm{Fl}}(M) \arrow[l,bend left=35,dashed,"\deleted{\unsure{F}}\unsure{\Phi|\mathcal F|F}_\alpha"]
\end{tikzcd}}
, $\Phi$ comme ``factorisation''
